% Supplementary Material for PAPER2_RECT_FREEFREE_DRAFT.tex
%   S.1  closed-contour four-corner Kirchhoff jump
%   S.2  in-plane conversion and Bardell (1996) mapping
%   S.3  full 92-row in-plane FE match list
%   S.4  cluster job-number provenance
%   S.5  flexural MAC against half-model SHELL281 eigenvectors
\documentclass[11pt]{article}
\usepackage[margin=1in,letterpaper]{geometry}
\usepackage{amsmath,amssymb,booktabs,array,xcolor}
\usepackage{textcomp}
\usepackage[hidelinks]{hyperref}
\usepackage{placeins}
\usepackage{longtable}
\renewcommand{\topfraction}{0.9}
\renewcommand{\bottomfraction}{0.85}
\renewcommand{\textfraction}{0.07}
\renewcommand{\floatpagefraction}{0.85}
\setcounter{topnumber}{4}
\setcounter{bottomnumber}{4}
\setcounter{totalnumber}{8}
\emergencystretch=2.5em

\renewcommand{\thesection}{S.\arabic{section}}
\renewcommand{\thesubsection}{S.\arabic{section}.\arabic{subsection}}
\renewcommand{\thetable}{S.\arabic{table}}
\renewcommand{\theequation}{S.\arabic{equation}}

\title{\bfseries Supplementary Material:\\
Exact Wave-Function Solution of Completely Free Rectangular Plates:
Four-Corner Kirchhoff Jump and Mode Screening}
\author{Garrek McCune$^{a,*}$, Daniel Stutts$^{a}$\\
\small\itshape $^a$Mechanical and Aerospace Engineering,\\
Missouri University of Science and Technology, Rolla, MO 65409, USA\\
\small\upshape $^*$Corresponding author. E-mail: ghmkfh@umsystem.edu}
\date{DRAFT --- \today}

\begin{document}
\maketitle

\noindent
This file is the complete Supplementary Material to the companion
manuscript: \S\S{}S.1--S.5. Main-text \S3.2 cites \S{}S.1; \S5.2--\S5.4 cite
\S\S{}S.2--S.3 (Gorman uses the same conversion as Bardell); the flexural
MAC of \S4.3 cites \S{}S.5; footnotes that read ``Job-level detail'' point
here at \S{}S.4.
Equation tags (16), (1) and (36) retain the source-paper numbering of
Seok, Tiersten and Scarton.

\section{Closed-contour four-corner Kirchhoff jump}\label{sec:sm-corner}

\subsection{The exact-differential term of Eq.~(16)}\label{sec:sm-eq16}
Seok, Tiersten and Scarton's rectangular Part~1 Eq.~(16) is a
closed-contour weak-form identity on the plate boundary
$\partial\Omega=c_C\cup c_N$ (clamped and natural portions). Written out
in full,
\begin{equation}
\begin{aligned}
&\int_{c_C}\Bigl\{u_3^{(0)}\,\delta\bigl(n_a t_{a3}^{(0)}\bigr)
  - u_{3,b}^{(0)}\,\delta\bigl(n_a t_{ab}^{(1)}\bigr)\Bigr\}\,ds \\
&\quad-\int_{c_N}\Bigl\{\bigl(n_a t_{a3}^{(0)}+(n_a t_{ab}^{(1)} s_b)_{,s}\bigr)\,
  \delta u_3^{(0)} - n_a t_{ab}^{(1)} n_b\,\delta u_{3,n}^{(0)}\Bigr\}\,ds \\
&\quad+\int_{c_N}\bigl(n_a t_{ab}^{(1)} s_b\,\delta u_3^{(0)}\bigr)_{,s}\,ds = 0 .
\end{aligned}
\label{eq:sm-16}
\end{equation}
Only the third integral is an exact differential in the arclength $s$.
Writing $\varphi := n_a t_{ab}^{(1)} s_b$ (the traction-moment contraction
that reduces to the twisting moment $M_{xy}$ on a smooth edge),
\begin{equation}
\int_{c_N}(\varphi\,\delta u_3^{(0)})_{,s}\,ds
= \sum_{\mathrm{corners}}\Bigl[\varphi\,\delta u_3^{(0)}\Bigr]_{\mathrm{end}}
  - \Bigl[\varphi\,\delta u_3^{(0)}\Bigr]_{\mathrm{start}},
\label{eq:sm-phi}
\end{equation}
the sum running over every free--free corner where two arcs of $c_N$
meet; interior-arc contributions cancel against the $\varphi_{,s}$ term
already present in the second integral (the usual Kirchhoff
effective-shear reduction). A completely free rectangle has
$c_C=\varnothing$ and four such corners.

Place one corner at the local origin with both positive local axes
pointing into the plate, so the two free edges meeting there are the
local $+x$ and $+y$ half-axes. Consistent counterclockwise orientation
of $\partial\Omega$ gives, on the vertical edge ($x=0^+$, $s$ increasing
in $+y$), $\mathbf{n}=+\mathbf{e}_x$, $\mathbf{s}=+\mathbf{e}_y$, so
$\varphi_V=t_{xy}^{(1)}$; on the horizontal edge ($y=0^+$, $s$ increasing
in $-x$), $\mathbf{n}=+\mathbf{e}_y$, $\mathbf{s}=-\mathbf{e}_x$, so
$\varphi_H=-t_{yx}^{(1)}$. Because $t_{ab}^{(1)}$ is symmetric,
$\varphi_V-\varphi_H=2\,t_{xy}^{(1)}$, so the corner's contribution to
the first variation is $+2\,t_{xy}^{(1)}\,\delta u_3^{(0)}$. In the
isotropic Kirchhoff limit $t_{xy}^{(1)}$ is the classical twisting
moment $M_{xy}$, so the natural condition is $2M_{xy}=0$ --- Timoshenko's
concentrated corner reaction $R=2M_{xy}$. This single-corner trace and its
classical-limit check were carried out and passed against that
pre-registered criterion before any four-corner assembly was attempted.

The overall sign of that contribution is a convention rather than a
result. Which edge of $\partial\Omega$ carries the $+$ label is fixed by
the global orientation chosen for Eq.~\eqref{eq:sm-16}, and reversing that
choice reverses every boundary term of the stationarity statement
together. Only the \emph{relative} sign between the corner term and the
surface terms it is added to is physical, and that relative sign is set by
the same $\pm 1$ orientation factor applied corner-by-corner in
\S\ref{sec:sm-four}. Fixing $\varphi$'s sign as the project's other
free--free derivations do (annular, variational and Green forms) writes
the corner contribution as
\begin{equation}
\text{corner contribution} = -2\,t_{xy}^{(1)}\,\delta u_3^{(0)}\Big|_{\mathrm{corner}}.
\label{eq:sm-single}
\end{equation}

\subsection{Four corners, and the discrete trial/test evaluation}\label{sec:sm-four}
In the discrete basis, the twisting-moment factor at a corner is the
trial function's mixed second derivative
$\partial^2 w/\partial x_1\partial x_2$ there, and the paired test factor
is the test function's value $w$ there. Direct differentiation confirms
that every corner term is
$M_{xy}^{\mathrm{trial}}(\mathrm{corner})\cdot
w^{\mathrm{test}}(\mathrm{corner})$, matching the already-validated
clamped-free tip-corner term. Define, for a given trial branch $\xi'$
and test branch $\xi''$,
\begin{equation}
\begin{aligned}
p_{1t}&=\left.\frac{\partial\psi_1'}{\partial x_2}\right|_{x_2=+\pi/2},&
p_{1b}&=\left.\frac{\partial\psi_1'}{\partial x_2}\right|_{x_2=-\pi/2},\\
q_{0t}&=\left.\psi_0''\right|_{x_2=+\pi/2},&
q_{0b}&=\left.\psi_0''\right|_{x_2=-\pi/2},
\end{aligned}
\label{eq:sm-pq}
\end{equation}
and, with $U_{t1},V_{t0}$ the $x_1=+L$ (tip) trial/test factors and
$U_{w1},V_{w0}$ the $x_1=-L$ (wall) factors, the four raw corner terms
\begin{equation}
\begin{aligned}
\mathrm{TT}&=p_{1t}U_{t1}q_{0t}V_{t0},&
\mathrm{TB}&=p_{1b}U_{t1}q_{0b}V_{t0},\\
\mathrm{WT}&=p_{1t}U_{w1}q_{0t}V_{w0},&
\mathrm{WB}&=p_{1b}U_{w1}q_{0b}V_{w0}.
\end{aligned}
\label{eq:sm-tt}
\end{equation}
TT and WB share the local orientation used above (outward normal along
the trial axis, $+2$ sense) and are themselves a diagonal pair; TB and WT
sit at the other diagonal, where the same closed-contour threading picks
up the opposite sense,
because the inward-axis Jacobian there is $-1$ --- the discrete
restatement of the classical alternating-sign corner-reaction pattern at
the four corners of a free rectangular plate.

\subsection{The parity identity, in both symmetry classes}\label{sec:sm-parity}
The half-model SYM/ANTI mirror conditions fix the phase of the $m=0$ and
$m=1$ $x_2$-profiles at $x_2=\pm\pi/2$. Direct evaluation of that phase
algebra gives:
\begin{itemize}\itemsep2pt
\item \textbf{SYM} ($m=1$, extra $\pi/2$ phase flips sign under
$x_2\to-x_2$; $m=0$ does not): $p_{1b}=-p_{1t}$ exactly,
$q_{0b}=q_{0t}$ exactly.
\item \textbf{ANTI} (roles of the two profiles swap): $p_{1b}=p_{1t}$
exactly, $q_{0b}=-q_{0t}$ exactly.
\end{itemize}
Both were confirmed numerically to working precision (the bracket
$p_{1t}q_{0t}+p_{1b}q_{0b}$ reproduces $0$ to ${\sim}10^{-27}$ at
25 decimal digits, independent of branch, $\Lambda$, $\ell/b$ or $\nu$).
\textbf{The same two relations force $\mathrm{TB}=-\mathrm{TT}$ and
$\mathrm{WB}=-\mathrm{WT}$ in both classes}, term by term rather than
only in the combined bracket: in SYM,
$\mathrm{TB}=p_{1b}U_{t1}q_{0b}V_{t0}=(-p_{1t})U_{t1}(q_{0t})V_{t0}=-\mathrm{TT}$,
and identically $\mathrm{WB}=-\mathrm{WT}$; in ANTI,
$\mathrm{TB}=p_{1t}U_{t1}(-q_{0t})V_{t0}=-\mathrm{TT}$ and again
$\mathrm{WB}=-\mathrm{WT}$.

\textbf{Consequence 1 --- why the original four-term same-sign sum is
identically zero.} The first free-free corner term ever assembled was the
naive $\mathrm{TT}+\mathrm{TB}+\mathrm{WT}+\mathrm{WB}$. Substituting
$\mathrm{TB}=-\mathrm{TT}$, $\mathrm{WB}=-\mathrm{WT}$ gives
$\mathrm{TT}-\mathrm{TT}+\mathrm{WT}-\mathrm{WT}\equiv 0$ in both
classes --- confirmed numerically to $1.6\times 10^{-24}$ (SYM) and
$2.8\times 10^{-23}$ (ANTI) on real assembled matrices, i.e.\ exactly
zero to precision, not merely small. That combination contributed
nothing to $\mathbf{K}$ regardless of $\Lambda$, $\ell/b$, or branch.

\textbf{Consequence 2 --- the surviving checkerboard is exactly twice the
single-point shorthand.} The deployed corner term flips the sign of the
diagonal pair instead of summing everything the same way:
\begin{equation}
\mathrm{corner}_{ff} = (T-\nu)(-2)\,(\mathrm{TT}-\mathrm{TB}-\mathrm{WT}+\mathrm{WB}).
\label{eq:sm-checkerboard}
\end{equation}
Substituting the same two relations: in SYM,
$\mathrm{TT}-\mathrm{TB}-\mathrm{WT}+\mathrm{WB}
=\mathrm{TT}-(-\mathrm{TT})-\mathrm{WT}+(-\mathrm{WT})
=2(\mathrm{TT}-\mathrm{WT})$; in ANTI the identical substitution gives
the identical result, $2(\mathrm{TT}-\mathrm{WT})$. So the checkerboard
bracket equals \textbf{twice} the earlier single-point candidate
$\mathrm{TT}-\mathrm{WT}$ --- in both symmetry classes, by the same
two-line substitution.

\subsection{Numerical checks}\label{sec:sm-checks}
\begin{itemize}\itemsep2pt
\item \textbf{Regression on the unmodified path.} The clamped-free branch
of the out-of-plane assembler is untouched by this term (it uses its own
corner, not the free-free corner). A diagnostic swap of the free-free
formula into the clamped-free path reproduces the parent
byte-for-byte ($\max|\mathbf{K}-\mathbf{K}_{\mathrm{parent}}|=0$ at the
validated Seok Part~1 Table~3 SYM benchmark, $\ell/b=1.5$,
$\Lambda=0.1551$). The
solver version identifier is unchanged; only the free-free corner
formula moved.
\item \textbf{Algebraic gate.} Before any FE comparison was trusted, a
pre-registered three-way comparison required: the naive sum genuinely
$\equiv 0$; the single-point and checkerboard candidates genuinely
nonzero; and
$\mathbf{K}_{\mathrm{checkerboard}}-\mathbf{K}_{\mathrm{naive}}
=2\,(\mathbf{K}_{\mathrm{single}}-\mathbf{K}_{\mathrm{naive}})$
to ${\sim}10^{-10}$ relative. All three passed.
\item \textbf{Independent finite-element cross-check.} At the production
basis, the checkerboard's nearest-to-FE dip beat both the naive-sum term
(which contributes nothing) and the single-point candidate at every
finite-element anchor available at the time (four points across
$\ell/b=1.5,2.5$, SYM and ANTI). Main-text \S4's five-aspect-ratio
table is a later, larger-sample confirmation of the same formula.
\end{itemize}

\subsection{What this section does not claim}\label{sec:sm-limits}
The single-corner Eq.~(16) trace is a first-principles derivation,
checked against the classical $R=2M_{xy}$ limit. The extension from one
corner to the specific four-corner checkerboard sign pattern is argued
from the closed-contour orientation logic plus the classical
alternating-corner-reaction pattern, and is \emph{validated} rather than
independently re-derived corner-by-corner from Eq.~(16) a second time ---
the same status the annular companion's corner term has relative to its
own Green-identity cross-check. The parity identity is a fully
first-principles algebraic fact about this project's specific
$(m=0,m=1)$ trial-function evaluations, exact in both classes. Nothing
here is a claim about the in-plane formulation, which has no Kirchhoff
jump term at all (Part~2 Eq.~(1)/(36); see main-text \S5).

\section{In-plane nondimensionalization and Bardell mapping}\label{sec:sm-ip}

\subsection{This project's
\texorpdfstring{$\bar\Omega$}{Omega-bar}}\label{sec:sm-ombar}
Seok, Tiersten and Scarton Part~2 Eqs.~(22)--(23):
$\bar\omega=(\pi/(2b))\sqrt{c_{66}/\rho}$, $c_{66}=E/(2(1+\nu))$,
$\bar\Omega=\omega/\bar\omega$. For the FE decks ($b=1\,\mathrm{m}$,
$E=210\,\mathrm{GPa}$, $\nu=0.30$, $\rho=7800\,\mathrm{kg\,m^{-3}}$)
this is $\bar\Omega=f[\mathrm{Hz}]/804.481$, independent of $\ell/b$.
Half-model PLANE183, mesh-conv max $0.004\%$.

\subsection{Bardell, Langley and Dunsdon (1996)}\label{sec:sm-bardell}
N.~S. Bardell, R.~S. Langley and J.~M. Dunsdon, \emph{J.\ Sound Vib.}
\textbf{191}(3) (1996) 459--467 \cite{bardell1996}. Their parameter is
$\Omega_B=\omega a/C_L$ with $C_L^2=E/[\rho(1-\nu^2)]$, frequencies
printed under Figure~1 (F--F--F--F) rather than tabulated. Poisson ratio
is not stated numerically. Their own Table~1, an exact S--S--S--S
convergence study, gives the square pair $(0,1)=(1,0)=1.859$, which equals
$\pi\sqrt{(1-\nu)/2}=1.8586$ at $\nu=0.30$.

Plates identify as $a=2\ell$, so $a/b=\ell/b$ and
\begin{equation}
\Omega_B=\bar\Omega\cdot\pi\cdot(\ell/b)\cdot\sqrt{(1-\nu)/2}.
\label{eq:sm-bardell}
\end{equation}
Figure~1 caption: ``(a) $a/b=1$; (b) $a/b=2$''. Panel (a) is drawn 2:1
with unique frequencies $\{1.954, 2.961, 3.267, 4.726, 4.784, 5.205\}$;
panel (b) is square with the repeated pair $2.472,2.472$. Body text
(p.~462) assigns repeats to $a/b=1$. Mapping follows the rendered
proportions, not the caption letters. The source's own Table~1 does not
have this swap.

Printed F--F--F--F first six, assigned that way:
\begin{itemize}\itemsep2pt
\item $a/b=1$ (square, Fig.~1(b)): $2.321$, $2.472$, $2.472$, $2.628$,
$2.987$, $3.452$.
\item $a/b=2$ (2:1, Fig.~1(a)): $1.954$, $2.961$, $3.267$, $4.726$,
$4.784$, $5.205$.
\end{itemize}
All twelve convert onto matched $\bar\Omega^{*}$ in main-text Table~4
(solver max miss $1.69\%$; independent FE max miss $0.019\%$).

\section{Full in-plane FE match list}\label{sec:sm-92}
Window $\bar\Omega\in[0.02,2.50]$, five aspect ratios, both parities.
Count $5+4+8+6+10+7+13+10+16+13=92$ is every non-rigid FE target in
that window. Table~\ref{tab:ip92} lists the unique MATCH after collapsing
double-dips onto the same FE partner. Production discovery used
$n_{\mathrm{cpair}}=3$ and a $3\%$ cut; sixteen remaining targets were
recovered at $n_{\mathrm{cpair}}=5$ in a $\pm 0.05$ window centred on
the FE value.

\begin{center}
\footnotesize
\begin{longtable}{@{}rllrrr@{}}
\caption{Unique in-plane MATCH against independent PLANE183 half-model FE, $\bar\Omega\in[0.02,2.50]$, all five aspect ratios and both parities. Job numbers in the last column: 2456002 is $n_{\mathrm{cpair}}=3$ production discovery; 2456013 is the $n_{\mathrm{cpair}}=5$ mop-up of the sixteen remaining targets.}
\label{tab:ip92}\\
\toprule
$\ell/b$ & class & $\bar\Omega^{*}$ & $\bar\Omega_{\mathrm{FE}}$ & miss & job \\
\midrule
\endfirsthead
\caption[]{Unique in-plane MATCH (continued).}\\
\toprule
$\ell/b$ & class & $\bar\Omega^{*}$ & $\bar\Omega_{\mathrm{FE}}$ & miss & job \\
\midrule
\endhead
\bottomrule
\endfoot
1.0 & SYM & 1.3300 & 1.32984 & 0.01\% & 2456002 \\
1.0 & SYM & 1.4100 & 1.41422 & 0.30\% & 2456002 \\
1.0 & SYM & 1.6100 & 1.60733 & 0.17\% & 2456002 \\
1.0 & SYM & 1.8600 & 1.85745 & 0.14\% & 2456002 \\
1.0 & SYM & 2.0000 & 2.00319 & 0.16\% & 2456002 \\
\addlinespace
1.0 & ANTI & 1.2400 & 1.24858 & 0.69\% & 2456002 \\
1.0 & ANTI & 1.3298 & 1.32984 & 0.00\% & 2456013 \\
1.0 & ANTI & 2.0000 & 2.00319 & 0.16\% & 2456002 \\
1.0 & ANTI & 2.3200 & 2.31523 & 0.21\% & 2456002 \\
\addlinespace
1.5 & SYM & 1.0500 & 1.04551 & 0.43\% & 2456002 \\
1.5 & SYM & 1.4100 & 1.41236 & 0.17\% & 2456002 \\
1.5 & SYM & 1.4243 & 1.42432 & 0.00\% & 2456013 \\
1.5 & SYM & 1.6300 & 1.62439 & 0.35\% & 2456002 \\
1.5 & SYM & 1.6900 & 1.69877 & 0.52\% & 2456002 \\
1.5 & SYM & 2.1100 & 2.10886 & 0.05\% & 2456002 \\
1.5 & SYM & 2.2200 & 2.22185 & 0.08\% & 2456002 \\
1.5 & SYM & 2.2800 & 2.28542 & 0.24\% & 2456002 \\
\addlinespace
1.5 & ANTI & 0.7800 & 0.78782 & 0.99\% & 2456002 \\
1.5 & ANTI & 1.0333 & 1.03335 & 0.00\% & 2456013 \\
1.5 & ANTI & 1.5712 & 1.57124 & 0.00\% & 2456013 \\
1.5 & ANTI & 1.6400 & 1.64765 & 0.46\% & 2456002 \\
1.5 & ANTI & 2.2025 & 2.20253 & 0.00\% & 2456013 \\
1.5 & ANTI & 2.2400 & 2.24065 & 0.03\% & 2456002 \\
\addlinespace
2.0 & SYM & 0.8100 & 0.79652 & 1.69\% & 2456002 \\
2.0 & SYM & 1.4000 & 1.40011 & 0.01\% & 2456002 \\
2.0 & SYM & 1.4142 & 1.41422 & 0.00\% & 2456013 \\
2.0 & SYM & 1.4400 & 1.44333 & 0.23\% & 2456002 \\
2.0 & SYM & 1.6500 & 1.65355 & 0.21\% & 2456002 \\
2.0 & SYM & 1.7500 & 1.77390 & 1.35\% & 2456002 \\
2.0 & SYM & 1.8100 & 1.81566 & 0.31\% & 2456002 \\
2.0 & SYM & 2.0000 & 2.00440 & 0.22\% & 2456002 \\
2.0 & SYM & 2.3000 & 2.29856 & 0.06\% & 2456002 \\
2.0 & SYM & 2.3300 & 2.33104 & 0.04\% & 2456002 \\
\addlinespace
2.0 & ANTI & 0.5200 & 0.52557 & 1.06\% & 2456002 \\
2.0 & ANTI & 0.8800 & 0.87891 & 0.12\% & 2456002 \\
2.0 & ANTI & 1.2715 & 1.27148 & 0.00\% & 2456013 \\
2.0 & ANTI & 1.2800 & 1.28703 & 0.55\% & 2456002 \\
2.0 & ANTI & 1.7400 & 1.73452 & 0.32\% & 2456002 \\
2.0 & ANTI & 1.8400 & 1.84467 & 0.25\% & 2456002 \\
2.0 & ANTI & 2.1400 & 2.13652 & 0.16\% & 2456002 \\
\addlinespace
2.5 & SYM & 0.6400 & 0.64049 & 0.08\% & 2456002 \\
2.5 & SYM & 1.2200 & 1.22293 & 0.24\% & 2456002 \\
2.5 & SYM & 1.4100 & 1.41225 & 0.16\% & 2456002 \\
2.5 & SYM & 1.4266 & 1.42658 & 0.00\% & 2456013 \\
2.5 & SYM & 1.5400 & 1.53449 & 0.36\% & 2456002 \\
2.5 & SYM & 1.6700 & 1.67003 & 0.00\% & 2456002 \\
2.5 & SYM & 1.6900 & 1.69350 & 0.21\% & 2456002 \\
2.5 & SYM & 1.7800 & 1.78140 & 0.08\% & 2456002 \\
2.5 & SYM & 2.0500 & 2.04232 & 0.38\% & 2456002 \\
2.5 & SYM & 2.0700 & 2.07671 & 0.32\% & 2456002 \\
2.5 & SYM & 2.1300 & 2.13048 & 0.02\% & 2456002 \\
2.5 & SYM & 2.3900 & 2.38996 & 0.00\% & 2456002 \\
2.5 & SYM & 2.4047 & 2.40468 & 0.00\% & 2456013 \\
\addlinespace
2.5 & ANTI & 0.3800 & 0.37585 & 1.10\% & 2456002 \\
2.5 & ANTI & 0.7200 & 0.71776 & 0.31\% & 2456002 \\
2.5 & ANTI & 1.0600 & 1.06883 & 0.83\% & 2456002 \\
2.5 & ANTI & 1.1125 & 1.11248 & 0.00\% & 2456013 \\
2.5 & ANTI & 1.4400 & 1.44246 & 0.17\% & 2456002 \\
2.5 & ANTI & 1.5400 & 1.53132 & 0.57\% & 2456002 \\
2.5 & ANTI & 1.8000 & 1.80308 & 0.17\% & 2456002 \\
2.5 & ANTI & 2.0000 & 1.99255 & 0.37\% & 2456002 \\
2.5 & ANTI & 2.2000 & 2.14745 & 2.45\% & 2456002 \\
2.5 & ANTI & 2.2800 & 2.34878 & 2.93\% & 2456002 \\
\addlinespace
3.0 & SYM & 0.5400 & 0.53502 & 0.93\% & 2456002 \\
3.0 & SYM & 1.0600 & 1.04554 & 1.38\% & 2456002 \\
3.0 & SYM & 1.4142 & 1.41422 & 0.00\% & 2456013 \\
3.0 & SYM & 1.4199 & 1.41791 & 0.14\% & 2456013 \\
3.0 & SYM & 1.4200 & 1.41978 & 0.02\% & 2456002 \\
3.0 & SYM & 1.6000 & 1.60066 & 0.04\% & 2456002 \\
3.0 & SYM & 1.6635 & 1.66346 & 0.00\% & 2456013 \\
3.0 & SYM & 1.6900 & 1.67627 & 0.82\% & 2456002 \\
3.0 & SYM & 1.8000 & 1.80401 & 0.22\% & 2456002 \\
3.0 & SYM & 1.8400 & 1.83870 & 0.07\% & 2456002 \\
3.0 & SYM & 2.0000 & 2.00478 & 0.24\% & 2456002 \\
3.0 & SYM & 2.1600 & 2.16287 & 0.13\% & 2456002 \\
3.0 & SYM & 2.2100 & 2.21145 & 0.07\% & 2456002 \\
3.0 & SYM & 2.2700 & 2.27554 & 0.24\% & 2456002 \\
3.0 & SYM & 2.4300 & 2.44758 & 0.72\% & 2456002 \\
3.0 & SYM & 2.4600 & 2.46034 & 0.01\% & 2456002 \\
\addlinespace
3.0 & ANTI & 0.2800 & 0.28169 & 0.60\% & 2456002 \\
3.0 & ANTI & 0.5800 & 0.57789 & 0.37\% & 2456002 \\
3.0 & ANTI & 0.8800 & 0.88739 & 0.83\% & 2456002 \\
3.0 & ANTI & 1.0315 & 1.03153 & 0.00\% & 2456013 \\
3.0 & ANTI & 1.2702 & 1.27024 & 0.00\% & 2456013 \\
3.0 & ANTI & 1.3000 & 1.29465 & 0.41\% & 2456002 \\
3.0 & ANTI & 1.5800 & 1.57876 & 0.08\% & 2456002 \\
3.0 & ANTI & 1.6600 & 1.66509 & 0.31\% & 2456002 \\
3.0 & ANTI & 1.9000 & 1.91234 & 0.65\% & 2456002 \\
3.0 & ANTI & 2.0000 & 1.99856 & 0.07\% & 2456002 \\
3.0 & ANTI & 2.2400 & 2.25049 & 0.47\% & 2456002 \\
3.0 & ANTI & 2.2800 & 2.27620 & 0.17\% & 2456002 \\
3.0 & ANTI & 2.4937 & 2.49371 & 0.00\% & 2456013 \\
\end{longtable}
\end{center}


\section{Cluster job-number provenance}\label{sec:provenance}
The main text states every result in prose, with no job number in its own
running text; this section is the audit trail. Each entry is keyed to the
main-text location that cites this section.

\begin{itemize}\itemsep2pt
\item \textbf{\S3.2 / \S{}S.1, checkerboard corner.} Deployed in the
out-of-plane assembler under the free-free option only. Algebraic gate
and four-point FE cross-check:
\path{probe_rect_ff_oop_derived_vs_fix_2026-09-03.py}. Clamped-free
regression at $\ell/b=1.5$, $\Lambda=0.1551$:
$\max|\mathbf{K}-\mathbf{K}_{\mathrm{parent}}|=0$.
\texttt{SOLVER\_VERSION} \texttt{2026-07-10.s10}, unbumped.
\item \textbf{\S3.3, Screen~B freeze.} Persist $n_{\mathrm{cpair}}=3$,
$|\Delta\Lambda|\le 0.01$, frozen before the extra-$\ell/b$ geometries
were run. Not retuned against FE. Post-hoc continuum sensitivity check:
job 2457815 rescored all 148 published flexural and in-plane matches
against the frozen fixed-step scan (all pass); job 2458006 re-scanned
the three that sat at the scan's outer edge with a $4\times$ finer grid;
job 2458121 settled them with a golden-section continuum minimisation.
Two main-text Table~2 entries ($\ell/b=1.5$ SYM $5.380$,
$\ell/b=2.5$ ANTI $6.260$) exceed the frozen cut at the continuum
optimum ($|\Delta\Lambda|=0.0103$, $0.0110$); one in-plane entry
($\ell/b=3.0$ ANTI $\bar\Omega^{*}=2.2400$, Table~\ref{tab:ip92}) does
not ($|\Delta\bar\Omega|=0.0094$).
\item \textbf{\S4, flexural FE.} $\ell/b=1.5$ half-model, job 2398617
lineage; $\ell/b=1.0,2.5$, P5 Phase~A queue; $\ell/b=2.0,3.0$, job
2454654. Primary refine $1.0/1.5/2.5$: job 2453874 (12 unique MATCH at
$3\%$); $2.0/3.0$: job 2454701 (12 unique MATCH at $3\%$). Persist-only
SYM hits at $0.543$ and $0.670$: job 2454702, both HIT at $0.00\%$.
High-$\Lambda$ $1.0/1.5/2.5$: job 2453875; $2.0/3.0$: job 2454703
(machine zeros tagged \texttt{ZERO\_ART}).
\item \textbf{\S4.4, $\Lambda_{\mathrm{FE}}=1.982$ unmatched
curiosity.} Job 2455675 reopened $\ell/b=1.0$ SYM $1.98249$ at persist
basis $(6,3)$, step $0.0002$ in $[1.96249,2.00249]$: global min
$\sigma=1.702\times 10^{-6}$ at the window's right edge, not flanked
(pre-registered reopen test failed; miss $1.009\%$). Job 2456003 widened
the scan to $[1.73,2.23]$ (step $0.002$): $\sigma$ stays within one
order of magnitude of its floor (ratio $3.12$), no dip within $3\%$ of
$1.98249$. Flat-plateau confirmed; Screen~B not retuned.
\item \textbf{\S4.5, Leissa 1973.} Conversion of Table~C15
$\lambda=\omega a^2\sqrt{\rho h/D}$ through the same
$K=24.6644\,\mathrm{Hz}$ as \S4; no cluster job. Fifteen overlapping
entries, each paired with its own parity under Leissa's $(m,n)$ rather
than with the nearest frequency; three named near-misses (square
$\lambda=19.789$; $\ell/b=1.5$ $\lambda=58.201$ and $\lambda=67.494$)
discussed in the main text.
\item \textbf{\S5.1, in-plane no-corner.} Seok Part~2 Eq.~(1)/(36);
assembler is wall$+$tip only. Free-free option is an edge-list swap.
Default clamped-free bit-identical (sandbox
$\max|\mathbf{K}_{\mathrm{default}}-\mathbf{K}_{\mathrm{cf}}|=0$).
\item \textbf{\S5.2, in-plane FE.} Job 2455569, PLANE183 half-model at
five aspect ratios, $5/5$ OK, mesh-conv max $0.004\%$,
$\bar\Omega=f/804.481$.
\item \textbf{\S5.2, in-plane discovery.} Job 2455570
($n_{\mathrm{cpair}}=0$, $5\%$ cut): 38 unique MATCH, ANTI branch weak.
Jobs 2455676/77/78 traced this to basis starvation, not missing physics.
Job 2456002 (production re-discovery at $n_{\mathrm{cpair}}=3$, $3\%$
cut): 76 unique MATCH. Job 2456013 (targeted mop-up,
$n_{\mathrm{cpair}}=5$, $\pm 0.05$ on the remaining 16 FE values):
$16/16$ recovered, $0.00$--$0.14\%$ miss. Net $92/92$ in
$\bar\Omega\in[0.02,2.50]$. Two apparent losses between the 38- and
76-match passes, $(2.0,\mathrm{ANTI},1.27148)$ and
$(3.0,\mathrm{SYM},1.66346)$, are both present in the 92
(root-crowding). Job 2455677's three ANTI ``NO DIP'' results
($\ell/b=1.0$ ANTI) are the same seeding artifact: their true FE
targets $1.24858/2.00319/2.31523$ are clean matches in the 92.
\item \textbf{\S5.3, Bardell 1996.} Frequencies read from Figure~1 of
the source PDF (mode-plot labels, not a table). Aspect-ratio mapping by
plate proportions and repeated-frequency physics, not the caption
letters. Conversion Eq.~\eqref{eq:sm-bardell} at $\nu=0.30$, confirmed
by the source's Table~1 S--S--S--S exact $1.859$. Matched
$\bar\Omega^{*}$ from jobs
2456002 and 2456013; FE from job 2455569. $12/12$ published F--F--F--F
first-six values matched.
\item \textbf{\S5.4, Gorman 2004.} Tabulated $\lambda^2$ doubled before
Eq.~\eqref{eq:sm-bardell} (Gorman's $a$ is the half-length); nineteen
in-window eigenvalues, no cluster job.
\item \textbf{\S4.3, flexural MAC.} Eigenvectors from job 2456738 (five
SHELL281 MAC-extract decks, $5/5$ \texttt{QUEUE\_OK}). Convert
$f_{\mathrm{Hz}}/24.6644$ in Python; the deck-printed $\Lambda$ column
is Instance~3 of the project's \texttt{*VWRITE}/KFAC bug and is not
used. Off-cluster reconstruction:
\path{probe_rect_ff_oop_mac_2026-09-05.py} and
\path{probe_rect_ff_oop_mac_strengthen_2026-09-05.py}. Bar $0.744$
is the annular companion's production call, unre-tuned. Job 2456779
(dps$=40$ persist golden section on the seven main-text Table~1 ANTI rows
that do not meet the bar): interior minima within $8\times 10^{-4}$ of the
published
$\Lambda^{*}$; persist MAC against the frequency partner does not
cross $0.744$ (gate~A of that probe; the bar is not retuned).
\texttt{SOLVER\_VERSION} \texttt{2026-07-10.s10}, unbumped.
\end{itemize}

\section{Flexural MAC against half-model eigenvectors}\label{sec:sm-mac}
The identity-weighted MAC is
$\mathrm{MAC}=|\langle u_Z,w\rangle|^2/(\|u_Z\|^2\|w\|^2)$ on the
phase-aligned real part of the reconstructed analytical $w$, evaluated
at the FE nodes of the half-model ($X\in[-\ell,\ell]$, $Y\in[0,b]$,
$x_1=X\pi/(2b)$, $x_2=Y\pi/(2b)$). The bar $0.744$ is the annular companion's production call
\cite{mccune2026annular}, not retuned against these data.
Table~\ref{tab:oop-mac}
lists every main-text Tables~1--2 pair that meets the bar against its
frequency-matched same-parity FE mode: $36$ of the $56$ published rows,
$25$ of $29$ symmetric and $11$ of $27$ antisymmetric. Included are three
persist-basis recoveries, six $n_{\mathrm{cpair}}=6$ recoveries (three of those after
discarding a machine-zero singular vector at $\sigma\sim 10^{-14}$,
the known $\Lambda=4,6$ artifact), and the persist-only
$\ell/b=3.0$ SYM $0.6698$ row whose persist kernel is two-dimensional
(the next singular vector is the physical shape). Known Screen~B leaks
are not in the table: SYM $0.41$ has MAC${}\le 0.05$ against every FE
mode of that parity; ANTI $0.46$ stays below the bar. Rows that remain
below the bar --- twenty rows, including the $\ell/b=2.5$ first ANTI (MAC
saturates at $0.655$), the $\ell/b=1.5$ ANTI $2.124$ weakest primary, and
both $^\ddagger$-marked rows of main-text Table~2 --- are frequency
matches only. The lowest MAC in the table is $0.751$ and the highest value
reached by any excluded row is $0.655$, so the confirmed set is unchanged
for any bar between those two.

% Confirmed MAC rows for Paper 2 SM S.5. Generated from
% probe_rect_ff_oop_mac_2026-09-05.py + n_cpair=6 strengthen pass.
% Bar 0.744 unre-tuned. Do not add SHAPE_MISMATCH rows here.
\begin{longtable}{@{}rllrrrl@{}}
\caption{MAC of analytical $w$ against the frequency-matched SHELL281
$U_Z$. Production basis unless noted. Bar $0.744$, not retuned.}\\
\label{tab:oop-mac}\\
\toprule
$\ell/b$ & class & $\Lambda^{*}$ & $\Lambda_{\mathrm{FE}}$ & MAC & basis & note \\
\midrule
\endfirsthead
\toprule
$\ell/b$ & class & $\Lambda^{*}$ & $\Lambda_{\mathrm{FE}}$ & MAC & basis & note \\
\midrule
\endhead
1.0 & ANTI & 1.366 & 1.35288 & 0.941 & prod. & 1st ANTI; Leissa $(2,2)$ \\
1.5 & SYM  & 0.964 & 0.96347 & 1.000 & prod. & 1st SYM; Leissa $(3,1)$ \\
1.5 & SYM  & 2.254 & 2.24373 & 1.000 & prod. & Leissa $(1,3)$ \\
1.5 & ANTI & 0.906 & 0.89838 & 0.852 & prod. & 1st ANTI; Leissa $(2,2)$ \\
2.0 & SYM  & 0.5434 & 0.54338 & 1.000 & persist & persist-only frequency \\
2.0 & SYM  & 1.510 & 1.50735 & 0.849 & persist & \\
2.0 & SYM  & 2.234 & 2.22558 & 0.999 & prod. & \\
2.0 & ANTI & 0.674 & 0.66858 & 0.751 & prod. & 1st ANTI \\
2.5 & SYM  & 0.348 & 0.34766 & 0.999 & prod. & 1st SYM; Leissa $(3,1)$ \\
2.5 & SYM  & 0.978 & 0.96559 & 0.871 & persist & \\
2.5 & SYM  & 1.888 & 1.88370 & 0.999 & prod. & \\
2.5 & SYM  & 2.278 & 2.27059 & 0.999 & prod. & \\
2.5 & SYM  & 3.190 & 3.19522 & 0.915 & $n_{\mathrm{cpair}}=6$ & close pair; see text \\
2.5 & SYM  & 4.180 & 4.13819 & 0.956 & $n_{\mathrm{cpair}}=6$ & skip machine-zero SV \\
2.5 & ANTI & 1.918 & 1.90392 & 1.000 & prod. & Leissa $(4,2)$ \\
2.5 & ANTI & 6.400 & 6.36131 & 1.000 & prod. & \\
3.0 & SYM  & 0.242 & 0.24127 & 0.959 & prod. & 1st SYM \\
3.0 & SYM  & 0.6698 & 0.66984 & 1.000 & persist & next SV of a 2-D kernel \\
3.0 & SYM  & 1.320 & 1.31785 & 1.000 & prod. & \\
3.0 & SYM  & 2.162 & 2.15417 & 0.779 & persist & \\
3.0 & SYM  & 2.256 & 2.24841 & 0.999 & prod. & \\
3.0 & SYM  & 2.480 & 2.45935 & 0.993 & $n_{\mathrm{cpair}}=6$ & \\
3.0 & SYM  & 3.340 & 3.32528 & 0.993 & prod. & \\
3.0 & SYM  & 3.660 & 3.62467 & 0.964 & $n_{\mathrm{cpair}}=6$ & skip machine-zero SV \\
3.0 & SYM  & 5.660 & 5.59746 & 0.926 & $n_{\mathrm{cpair}}=6$ & skip machine-zero SV \\
3.0 & ANTI & 0.444 & 0.44044 & 1.000 & prod. & 1st ANTI; $0.46$-family \\
3.0 & ANTI & 3.160 & 3.14253 & 0.999 & prod. & \\
3.0 & ANTI & 5.560 & 5.51872 & 0.963 & prod. & \\
3.0 & ANTI & 6.380 & 6.34339 & 1.000 & prod. & \\
1.0 & SYM  & 2.460 & 2.45439 & 0.997 & prod. & Leissa $(3,1)$ \\
1.0 & SYM  & 6.190 & 6.16026 & 0.867 & $n_{\mathrm{cpair}}=6$ & \\
1.0 & SYM  & 6.460 & 6.36553 & 1.000 & prod. & \\
1.5 & ANTI & 3.860 & 3.83354 & 0.993 & prod. & \\
2.0 & SYM  & 3.000 & 2.99812 & 0.862 & prod. & \\
2.0 & SYM  & 3.650 & 3.62227 & 1.000 & prod. & \\
2.0 & ANTI & 2.580 & 2.55180 & 0.999 & prod. & \\
\bottomrule
\end{longtable}


\begin{thebibliography}{99}\small
\setlength{\itemsep}{0pt}
\setlength{\parsep}{0pt}
\bibitem{seok3} J. Seok, H.F. Tiersten, H.A. Scarton, ``Free vibrations of
rectangular cantilever plates. Part 1: out-of-plane motion,'' J. Sound
Vib. 271 (2004) 131--146. \url{https://doi.org/10.1016/S0022-460X(03)00365-1}.
\bibitem{seok4} J. Seok, H.F. Tiersten, H.A. Scarton, ``Free vibrations of
rectangular cantilever plates. Part 2: in-plane motion,'' J. Sound Vib.
271 (2004) 147--158. \url{https://doi.org/10.1016/S0022-460X(03)00366-3}.
\bibitem{bardell1996} N.S. Bardell, R.S. Langley, J.M. Dunsdon, ``On the
free in-plane vibration of isotropic rectangular plates,'' J. Sound Vib.
191 (3) (1996) 459--467. \url{https://doi.org/10.1006/jsvi.1996.0134}.
\bibitem{mccune2026annular} G. McCune, D. Stutts, ``Exact wave-function
solution and spurious-mode discrimination for completely free annular
sector plates,'' companion paper, submitted simultaneously to this
journal (2026).
\end{thebibliography}
\end{document}
